\chapter{Considerações finais}

Este trabalho caracterizou a adequação prática e os limites de um \textit{pipeline} avaliativo de reconstrução de cenários 3D fotorrealistas para VR \textit{standalone}, articulando NeRF e \textit{3D Gaussian Splatting}. A contribuição não consiste em propor um novo método de reconstrução, mas em documentar um caminho reprodutível de treinamento, exportação, integração em Unity com UnitySplats e execução nativa no Meta Quest 3S, acompanhado de métricas e evidências de limitações.

No estudo de caso \textit{ns\_poster}, o NeRF foi utilizado como referência de fidelidade visual e o 3D-GS como representação voltada à renderização interativa. A comparação entre variantes mostrou que reduzir deterministicamente o número de gaussianas melhora o desempenho no headset, mas impõe perda mensurável de fidelidade offline e não elimina problemas de estabilidade visual. A poda 100k obteve PSNR 22,716, SSIM 0,80851 e LPIPS 0,31945; a 50k obteve 17,514, 0,60303 e 0,51890, respectivamente, ante 32,989, 0,95363 e 0,10495 da \textit{baseline}. A configuração de 50 mil gaussianas alcançou 58,37 FPS na série de cinco execuções, enquanto a \textit{baseline} alcançou 22,20 FPS. Nenhuma das variantes comparativas atingiu a referência de 90 FPS definida no protocolo.

O diagnóstico da variante \textit{splatfacto-big} reforçou esse compromisso. Ela apresentou indícios qualitativos de maior fidelidade na condição observada, mas teve 15,24 FPS em uma execução estática e não foi considerada adequada para desempenho \textit{standalone}. Esses resultados sustentam que qualidade offline ou inspeção visual favorável não são evidência suficiente de viabilidade em VR: a avaliação precisa incluir o custo da representação, o tempo de quadro e a estabilidade da execução nativa.

As conclusões são delimitadas por uma única cena, pose fixa, cinco execuções somente para a série principal e avaliação das podas por substituição pós-exportação dos parâmetros Gaussianos. A observação visual estruturada é qualitativa, e os relatos de desconforto durante movimento de cabeça podem ter sido influenciados pela sucessão dos testes. A observação visual do diagnóstico de maior capacidade também ocorreu em uma única execução e não constitui medida perceptual. Portanto, os resultados não devem ser generalizados para qualquer ambiente, dispositivo ou configuração de captura. O trabalho permanece restrito a cenas estáticas e não contempla reconstrução dinâmica, \textit{relighting} avançado, áudio espacial ou interação física complexa.

Dessa forma, o protótipo cumpre seu papel avaliativo: preserva tanto a melhoria de desempenho obtida pela poda quanto a degradação quantitativa e qualitativa que a acompanhou. Como continuidade, recomenda-se ampliar poses e cenas de teste e investigar técnicas de redução de custo que preservem melhor a estabilidade visual. Essas ações permitiriam avaliar se a aproximação da referência de desempenho pode ocorrer sem perda de fidelidade incompatível com a visualização imersiva.
